\documentclass[12pt,a4paper]{article}
\usepackage[utf8]{inputenc}
\usepackage[T1]{fontenc}
\usepackage{geometry}
\usepackage{enumitem}
\usepackage{fancyhdr}
\usepackage{lastpage}
\usepackage{hyperref}

\geometry{margin=2.5cm}
\pagestyle{fancy}
\fancyhf{}
\rhead{Page \thepage\ of \pageref{LastPage}}
\lhead{SCU101 - Quiz (en)}
\renewcommand{\headrulewidth}{0.4pt}

\title{Quiz: SCU101}
\author{Language: en}
\date{\today}

\begin{document}

\maketitle
\thispagestyle{fancy}

\section*{Instructions}
Answer all questions by selecting the best option (A, B, C, or D). The answer key is provided at the end of this document.

\bigskip


\subsection*{Question 1}
What is the main difference between private browsing and using the TOR network?

\begin{enumerate}[label=\Alph*., leftmargin=*]
\item TOR leaves local traces on your computer, while private browsing does not.
\item Tor masks the user's IP address, while Private browsing does not.
\item Private browsing masks the user's IP address, while TOR does not.
\item Private browsing and TOR are the same thing.
\end{enumerate}

\bigskip


\subsection*{Question 2}
Why is it recommended to limit the number of extensions installed on your browser?

\begin{enumerate}[label=\Alph*., leftmargin=*]
\item To prevent the browser from crashing.
\item To avoid potential security and performance issues.
\item To avoid being tracked by third-party sites.
\item To save storage space on your device.
\end{enumerate}

\bigskip


\subsection*{Question 3}
Why is it recommended to download software from the official website of the publisher rather than generic sites?

\begin{enumerate}[label=\Alph*., leftmargin=*]
\item To ensure the software is free from malicious software.
\item Because the software is usually cheaper on the official website.
\item To support the publisher financially.
\item Because the software on generic sites is usually outdated.
\end{enumerate}

\bigskip


\subsection*{Question 4}
What is the main advantage of using open-source software?

\begin{enumerate}[label=\Alph*., leftmargin=*]
\item It is always safer than closed-source software.
\item It is always more user-friendly than closed-source software.
\item It allows for verification that there is no hidden access to steal your personal data.
\item It is always free.
\end{enumerate}

\bigskip


\subsection*{Question 5}
Why is it recommended to delete cookies after each visit to a site?

\begin{enumerate}[label=\Alph*., leftmargin=*]
\item To prevent the site from showing you targeted ads.
\item To prevent the site from remembering your login information.
\item To prevent third-party sites from tracking your online activity.
\item To save storage space on your device.
\end{enumerate}

\bigskip


\subsection*{Question 6}
Why is it recommended to avoid the Google Chrome browser?

\begin{enumerate}[label=\Alph*., leftmargin=*]
\item Because it does not support extensions.
\item Due to its trackers.
\item Because it is not secure.
\item Because it is not user-friendly.
\end{enumerate}

\bigskip


\subsection*{Question 7}
What are cookies in the context of online browsing?

\begin{enumerate}[label=\Alph*., leftmargin=*]
\item Files created by websites to store information on your device.
\item Small pieces of data sent from a website and stored on the user's computer by the user's web browser.
\item Viruses that can harm your device.
\item Programs that track your online activity.
\end{enumerate}

\bigskip


\subsection*{Question 8}
What is private browsing?

\begin{enumerate}[label=\Alph*., leftmargin=*]
\item A feature that prevents internet browsers from storing browsing history, temporary internet files, form data, cookies, and user names and passwords.
\item A feature that hides your online activity from your internet service provider.
\item A feature that protects your device from viruses and other malicious software.
\item A feature that allows you to browse the internet anonymously.
\end{enumerate}

\bigskip


\subsection*{Question 9}
What is the TOR network?

\begin{enumerate}[label=\Alph*., leftmargin=*]
\item A network that protects your device from viruses and other malicious software.
\item A network that hides your online activity from your internet service provider.
\item A network that allows you to browse the internet faster.
\item A network that offers anonymity by masking the user's IP address and allowing access to the Darknet.
\end{enumerate}

\bigskip


\subsection*{Question 10}
What is the primary function of a VPN in a professional context?

\begin{enumerate}[label=\Alph*., leftmargin=*]
\item It allows employees to access the company's internal network remotely and securely.
\item It allows employees to completely block online advertising.
\item It allows employees to safely bypass geographical restrictions.
\item It allows employees to use public Wi-Fi networks securely.
\end{enumerate}

\bigskip


\subsection*{Question 11}
What is the difference between enterprise VPNs and consumer VPNs?

\begin{enumerate}[label=\Alph*., leftmargin=*]
\item Enterprise VPNs are usually more secure than consumer VPNs, which are always more expensive.
\item Consumer VPNs are always more expensive than enterprise VPNs due to the economic scale that the latter can access.
\item Due to their more exigent target audience, enterprise VPNs always encrypt data, while consumer VPNs don't.
\item Enterprise VPNs are more expensive and complex, while consumer VPNs are more accessible and user-friendly.
\end{enumerate}

\bigskip


\subsection*{Question 12}
What is the purpose of the HTTPS protocol?

\begin{enumerate}[label=\Alph*., leftmargin=*]
\item It encrypts data on websites to ensure secure exchange between the user and the website.
\item It allows users to access the internet without using private credentials.
\item It allows users to bypass geographical restrictions when paying via debit or credit card.
\item It allows users to block online advertising partially.
\end{enumerate}

\bigskip


\subsection*{Question 13}
Why is verifying the site's public key using a certificate system important?

\begin{enumerate}[label=\Alph*., leftmargin=*]
\item To ensure no scam can happen, but only when connected via a public Wi-Fi network.
\item To ensure that no malicious entity is impersonating the website.
\item To confirm the authenticity of the site.
\item To partially block online advertisements.
\end{enumerate}

\bigskip


\subsection*{Question 14}
Why is using European Wi-Fi access point providers, such as SNCF, safer?

\begin{enumerate}[label=\Alph*., leftmargin=*]
\item They provide faster internet speeds than American providers.
\item They block online advertisements if you have been identified as an European citizen.
\item They often allow users to bypass geographical restrictions.
\item They are not supposed to resell user connection data, due to GDPR regulations.
\end{enumerate}

\bigskip


\subsection*{Question 15}
What is one way to avoid online scams?

\begin{enumerate}[label=\Alph*., leftmargin=*]
\item Use public Wi-Fi networks.
\item Use sites that do not use the HTTPS protocol.
\item Use sites that do not have a padlock displayed.
\item Verify the identity of the site you are browsing, especially by checking the extension and domain name.
\end{enumerate}

\bigskip


\subsection*{Question 16}
Why is it important to prioritize reliability and technicality over popularity when choosing a VPN?

\begin{enumerate}[label=\Alph*., leftmargin=*]
\item Popular VPN providers are more likely to have geographical restrictions.
\item VPN providers that collect the least personal information are generally the safest.
\item Popular VPN providers are more likely to be targeted by hackers.
\item Popular VPN providers are more likely to have slower internet speeds.
\end{enumerate}

\bigskip


\subsection*{Question 17}
What is one benefit of using a VPN for browsing the internet at home?

\begin{enumerate}[label=\Alph*., leftmargin=*]
\item It saves mobile data consumption compared to not having an established VPN connection.
\item It enhances security for online traffic by reducing the threat posed by third parties.
\item It can increase the speed of your internet connection.
\item It always allows access to more websites.
\end{enumerate}

\bigskip


\subsection*{Question 18}
What is one way to confirm that you are on the site you intend to visit?

\begin{enumerate}[label=\Alph*., leftmargin=*]
\item Check the color scheme of the site.
\item Check the URLs and the small padlock in the address bar.
\item Check the speed of the site.
\item Check the number of advertisements on the site.
\end{enumerate}

\bigskip


\subsection*{Question 19}
What is the recommended practice for using an administrator account on Windows?

\begin{enumerate}[label=\Alph*., leftmargin=*]
\item Create two separate administrator accounts - one for daily use and one for operating system administration.
\item Use the administrator account daily if the PC is not connected to the internet.
\item Do not create an administrator account.
\item Create two separate accounts - an administrator account and an account for daily use.
\end{enumerate}

\bigskip


\subsection*{Question 20}
Why is it crucial to regularly update the operating system and applications?

\begin{enumerate}[label=\Alph*., leftmargin=*]
\item To slow down the machine.
\item To protect against new threats and vulnerabilities.
\item To make the machine obsolete.
\item To change the user interface.
\end{enumerate}

\bigskip


\subsection*{Question 21}
What is the built-in antivirus in Windows?

\begin{enumerate}[label=\Alph*., leftmargin=*]
\item Norton
\item McAfee
\item Windows Defender
\item Kaspersky
\end{enumerate}

\bigskip


\subsection*{Question 22}
What is the potential threat of using cracked software that cannot be updated?

\begin{enumerate}[label=\Alph*., leftmargin=*]
\item The software will stop working after a certain period.
\item The software will not have any new features.
\item The software will consume more system resources.
\item The arrival of a virus during its illegal download and an insecure use against new forms of attack.
\end{enumerate}

\bigskip


\subsection*{Question 23}
What is the recommended practice when inserting a USB key into your computer?

\begin{enumerate}[label=\Alph*., leftmargin=*]
\item Insert any USB key without checking.
\item Insert the USB key and then check for viruses.
\item Do not insert unknown or suspicious USB keys.
\item Always format the USB key before inserting.
\end{enumerate}

\bigskip


\subsection*{Question 24}
What is the importance of using a strong password for computer security?

\begin{enumerate}[label=\Alph*., leftmargin=*]
\item To prevent unauthorized access to the system.
\item To improve the user interface.
\item To make the system run faster.
\item To increase the storage capacity.
\end{enumerate}

\bigskip


\subsection*{Question 25}
What are some solutions for data encryption on different operating systems?

\begin{enumerate}[label=\Alph*., leftmargin=*]
\item BitLocker for Mac, LUKS for Windows, and the built-in option for Linux.
\item BitLocker for Linux, LUKS for Mac, and the built-in option for Windows.
\item LUKS for all operating systems.
\item BitLocker for Windows, LUKS for Linux, and the built-in option for Mac.
\end{enumerate}

\bigskip


\subsection*{Question 26}
What is the recommended practice for storing the password for data encryption?

\begin{enumerate}[label=\Alph*., leftmargin=*]
\item Store the password in a text file on the computer.
\item Share the password with trusted friends.
\item Write down the password on a piece of paper to keep it in a safe place.
\item Use a common password that is easy to remember.
\end{enumerate}

\bigskip


\subsection*{Question 27}
Do Linux and Mac often need antivirus?

\begin{enumerate}[label=\Alph*., leftmargin=*]
\item This is true
\item Only Mac needs antivirus.
\item Only Linux needs antivirus.
\item This is false
\end{enumerate}

\bigskip


\subsection*{Question 28}
What is the primary reason for not paying hackers for ransomware or other attacks?

\begin{enumerate}[label=\Alph*., leftmargin=*]
\item Paying the ransom can encourage cybercriminals to continue their malicious activities since the payment method is usually untraceable.
\item Paying the ransom does not guarantee the recovery of your data but can encourage hackers to attack other victims.
\item Paying the ransom is illegal and can encourage cybercriminals to continue their malicious activities.
\item Paying the ransom does not guarantee the recovery of your data and can encourage cybercriminals to continue their malicious activities.
\end{enumerate}

\bigskip


\subsection*{Question 29}
What is the first vector of attack for hackers?

\begin{enumerate}[label=\Alph*., leftmargin=*]
\item Online shopping websites
\item Emails
\item Banking websites
\item Social media
\end{enumerate}

\bigskip


\subsection*{Question 30}
What is a simple solution to protect your data from ransomware attacks?

\begin{enumerate}[label=\Alph*., leftmargin=*]
\item Regularly change your computer or reinstall your operating system (OS).
\item Install multiple antivirus software programs on your computer and perform periodic scans.
\item Regularly back up your important data on an external device or a secure online storage service.
\item You can use a public Wi-Fi network if you have nothing to hide on your device.
\end{enumerate}

\bigskip


\subsection*{Question 31}
What is the golden rule to protect yourself from fraudulent emails?

\begin{enumerate}[label=\Alph*., leftmargin=*]
\item Open all email attachments to check their content.
\item Carefully check the full name of the sender as well as the origin of the email.
\item Share your personal information only with known senders.
\item Always reply to suspicious emails.
\end{enumerate}

\bigskip


\subsection*{Question 32}
What should you do if you detect a virus on your computer?

\begin{enumerate}[label=\Alph*., leftmargin=*]
\item Disconnect it from the internet, perform a full antivirus scan, and delete infected files.
\item Delete all files on the computer.
\item Continue using the computer normally.
\item Turn off the computer and leave it off for a few hours.
\end{enumerate}

\bigskip


\subsection*{Question 33}
What is ransomware?

\begin{enumerate}[label=\Alph*., leftmargin=*]
\item A type of malicious software that encrypts user data and demands a ransom to decrypt it.
\item A type of antivirus software that is capable of cleaning computer cache and cookies upon the payment of a one-time fee.
\item A type of software used for data recovery, it is becoming quite common in Europe nowadays.
\item A type of software used for data backup, which usually performs one backup every 7 days.
\end{enumerate}

\bigskip


\subsection*{Question 34}
What is the potential downside of creating a cloud backup using services like ProtonMail Drive, Sync, or Google Drive?

\begin{enumerate}[label=\Alph*., leftmargin=*]
\item Your data is potentially on the internet and held by a trusted third party.
\item These services are potentially not reliable for data backup.
\item These services sometimes lose user data, and transferring it to another cloud provider can be difficult.
\item These services usually do not provide sufficient storage space and suffer from low connection speed.
\end{enumerate}

\bigskip


\subsection*{Question 35}
What malicious content can email attachments and images contain?

\begin{enumerate}[label=\Alph*., leftmargin=*]
\item Spam
\item Promotional content
\item Advertisements
\item Malware
\end{enumerate}

\bigskip


\subsection*{Question 36}
What is the purpose of phishing attempts?

\begin{enumerate}[label=\Alph*., leftmargin=*]
\item To invite you to events or webinars.
\item To extract sensitive information such as your credentials and passwords.
\item To send you spam emails.
\item To sell you products or services.
\end{enumerate}

\bigskip


\subsection*{Question 37}
What is the advantage of using a self-hosted password manager like Bitwarden or KeePass?

\begin{enumerate}[label=\Alph*., leftmargin=*]
\item They usually offer more features than third-party services, although they are more expensive.
\item You can control where your passwords are stored, offering more security and control.
\item They are easier and always more secure to use than third-party services.
\item They are cheaper and faster than third-party services.
\end{enumerate}

\bigskip


\subsection*{Question 38}
Why is it important to gradually replace your old passwords with those generated by a password manager?

\begin{enumerate}[label=\Alph*., leftmargin=*]
\item The password manager automatically deletes old passwords.
\item Password managers require all passwords to be generated by them.
\item Password managers can only store a limited number of passwords.
\item Passwords generated by a password manager are stronger and longer, offering more security.
\end{enumerate}

\bigskip


\subsection*{Question 39}
What is the main advantage of using a password manager?

\begin{enumerate}[label=\Alph*., leftmargin=*]
\item It automatically logs you into your online accounts.
\item It protects your computer from viruses.
\item It stores all your passwords in one place, making them easier to find.
\item It helps you create strong passwords and remember them, so you don't have to.
\end{enumerate}

\bigskip


\subsection*{Question 40}
What is the main disadvantage of using a self-hosted password manager?

\begin{enumerate}[label=\Alph*., leftmargin=*]
\item It requires more technical knowledge to use.
\item It is more expensive than third-party services.
\item It is not as secure as third-party services.
\item It does not offer as many features as third-party services.
\end{enumerate}

\bigskip


\subsection*{Question 41}
What is the purpose of implementing 2FA?

\begin{enumerate}[label=\Alph*., leftmargin=*]
\item To make it harder for you to access your own account
\item To track your online activities
\item To protect your online accounts against unauthorized access.
\item To increase the speed of your internet connection
\end{enumerate}

\bigskip


\subsection*{Question 42}
Which of the following is a good and secure method in 2FA authentication?

\begin{enumerate}[label=\Alph*., leftmargin=*]
\item A letter sent to the account holder's home address
\item A code generated by an application like Google Authenticator or Authy.
\item A phone call or an SMS code to the account holder
\item An email with a 10-minute expiring code inside it to the account holder
\end{enumerate}

\bigskip


\subsection*{Question 43}
Why is SMS not considered the best option for 2FA?

\begin{enumerate}[label=\Alph*., leftmargin=*]
\item It is too slow and expensive; a phone call will be more effective.
\item It only provides proof of possession of a phone number.
\item It's too complicated to set up on a technical level.
\item This method is vulnerable to attacks like the SIM swap scam.
\end{enumerate}

\bigskip


\subsection*{Question 44}
What is the advantage of one-time passwords (HOTP and TOTP) over SMS for 2FA?

\begin{enumerate}[label=\Alph*., leftmargin=*]
\item They are easier to remember and faster to type in.
\item They require cryptographic calculations and are stored in memory rather than locally.
\item They do not require cryptographic calculations and are stored locally and in memory.
\item They require cryptographic calculation and are stored locally rather than in memory.
\end{enumerate}

\bigskip


\subsection*{Question 45}
What is the golden rule of cybersecurity?

\begin{enumerate}[label=\Alph*., leftmargin=*]
\item Cybersecurity is a one-time setup
\item Cybersecurity is only for tech-savvy people
\item Cybersecurity is not necessary for personal accounts
\item Cybersecurity is a moving target that will adapt to your learning journey.
\end{enumerate}

\bigskip


\subsection*{Question 46}
Why are hardware tokens, such as USB keys or smart cards, considered to offer optimal security for 2FA?

\begin{enumerate}[label=\Alph*., leftmargin=*]
\item They generate only one password for every site and verify the URL before allowing the connection.
\item They generate a unique private key for each site and verify the URL before allowing the connection.
\item They are newer technology and expensive, so they must be better.
\item They are physically harder to steal, but if thieves succeed, they'll easily extract the secrets from the device.
\end{enumerate}

\bigskip


\subsection*{Question 47}
Why is biometric data considered less secure than the combination of knowledge and possession for 2FA?

\begin{enumerate}[label=\Alph*., leftmargin=*]
\item Biometric data is easier to extract even if kept offline on the authentication device.
\item Biometric data should remain on the authentication device and not be disclosed online.
\item Biometric data tools are more expensive to implement and therefore not cost effective compared to other solutions.
\item Biometric data is both harder to remember and use when using an authentication device.
\end{enumerate}

\bigskip


\subsection*{Question 48}
What are the recommended practices for optimal security with strong authentication?

\begin{enumerate}[label=\Alph*., leftmargin=*]
\item Use a secure email address and password manager, and adopt 2FA using YubiKeys.
\item Use a common password for all accounts, and avoid 2FA
\item Use a simple password and rely on biometric data for 2FA
\item Use a secure email address, but avoid password managers and 2FA
\end{enumerate}

\bigskip


\subsection*{Question 49}
What is a good starting point for a career in cybersecurity from an academic perspective?

\begin{enumerate}[label=\Alph*., leftmargin=*]
\item A degree in physical education
\item A degree in arts and humanities
\item A degree in business management
\item A solid education in computer science, information systems, or a related field
\end{enumerate}

\bigskip


\subsection*{Question 50}
Why is practical experience important in cybersecurity?

\begin{enumerate}[label=\Alph*., leftmargin=*]
\item It is more important than academic knowledge
\item It is the only way to learn about cybersecurity
\item It allows you to apply your theoretical knowledge in real-world situations
\item It is a requirement for all cybersecurity jobs
\end{enumerate}

\bigskip


\subsection*{Question 51}
What is one of the benefits of attending cybersecurity conferences and workshops?

\begin{enumerate}[label=\Alph*., leftmargin=*]
\item It guarantees a job in cybersecurity
\item It is the only way to learn about cybersecurity
\item It is a requirement for all cybersecurity jobs
\item It allows you to build connections with industry professionals
\end{enumerate}

\bigskip


\subsection*{Question 52}
What are some of the tools used in cybersecurity?

\begin{enumerate}[label=\Alph*., leftmargin=*]
\item AutoCAD, SolidWorks, Revit
\item Excel, PowerPoint, Word
\item Wireshark, Metasploit, Nmap
\item Photoshop, Illustrator, InDesign
\end{enumerate}

\bigskip


\subsection*{Question 53}
What are some of the skills needed in cybersecurity?

\begin{enumerate}[label=\Alph*., leftmargin=*]
\item Cooking, baking
\item Log analysis, incident response
\item Graphic design, video editing
\item Sales, marketing
\end{enumerate}

\bigskip


\subsection*{Question 54}
What is one of the aspects of cybersecurity that requires regular monitoring?

\begin{enumerate}[label=\Alph*., leftmargin=*]
\item The constant evolution of threats
\item The weather
\item The news
\item The stock market
\end{enumerate}

\bigskip


\subsection*{Question 55}
What are some of the programming languages used in cybersecurity?

\begin{enumerate}[label=\Alph*., leftmargin=*]
\item HTML, CSS, JavaScript
\item Swift, Kotlin, Ruby
\item PHP, Perl, Lua
\item Python, C, Java
\end{enumerate}

\bigskip


\subsection*{Question 56}
What aspects of network knowledge are needed in cybersecurity?

\begin{enumerate}[label=\Alph*., leftmargin=*]
\item POP3, IMAP, SSH
\item TCP/IP, VPN, firewall
\item DHCP, DNS, ARP
\item HTTP, FTP, SMTP
\end{enumerate}

\bigskip


\subsection*{Question 57}
What certification is recognized worldwide in the field of cybersecurity?

\begin{enumerate}[label=\Alph*., leftmargin=*]
\item CISSP
\item ITIL
\item PMP
\item SHA2
\end{enumerate}

\bigskip


\subsection*{Question 58}
Why is it not recommended to use browser extensions for automatic password filling?

\begin{enumerate}[label=\Alph*., leftmargin=*]
\item They can make the user more vulnerable to phishing attacks, increase the attack surface, and slow down browser performance.
\item They are difficult to use.
\item They are not compatible with all browsers.
\item They require a lot of storage space.
\end{enumerate}

\bigskip


\subsection*{Question 59}
What are the advantages of using Passkeys as an alternative to traditional passwords?

\begin{enumerate}[label=\Alph*., leftmargin=*]
\item They can be used on all websites without any updates.
\item They do not require any additional authentication factors.
\item They offer encrypted random keys combined with a local factor. The keys are hosted by a provider but remain out of their reach.
\item They are easier to remember than traditional passwords.
\end{enumerate}

\bigskip


\subsection*{Question 60}
What is the main concern with the Single Sign-On (SSO) offered by internet giants?

\begin{enumerate}[label=\Alph*., leftmargin=*]
\item It requires a lot of storage space.
\item It poses problems in terms of availability and risks of censorship.
\item It is difficult to use.
\item It is not compatible with all websites.
\end{enumerate}

\bigskip


\subsection*{Question 61}
What is a common cause of security vulnerabilities?

\begin{enumerate}[label=\Alph*., leftmargin=*]
\item Slow internet speed.
\item Power outages.
\item Human negligence.
\item Outdated hardware.
\end{enumerate}

\bigskip


\subsection*{Question 62}
What are the advantages of using DOH (DNS over HTTPS) and DOT (DNS over TLS)?

\begin{enumerate}[label=\Alph*., leftmargin=*]
\item They block all advertisements.
\item They allow access to blocked websites.
\item They speed up the internet connection.
\item They encrypt the DNS connection, providing greater privacy and security, and avoiding man-in-the-middle attacks.
\end{enumerate}

\bigskip


\subsection*{Question 63}
What is the purpose of Lightning authentication?

\begin{enumerate}[label=\Alph*., leftmargin=*]
\item It allows you to use the same password for all services.
\item Allows anyone, even without a Lightning wallet, to authenticate without an email address, if a valid ID is shown.
\item It generates a different identifier for each service, without providing an email address or personal information.
\item It sends reminders to change passwords every month.
\end{enumerate}

\bigskip


\subsection*{Question 64}
What is a recommended practice for managing the master password of a password manager?

\begin{enumerate}[label=\Alph*., leftmargin=*]
\item Use the same password for all accounts.
\item Store it in your email.
\item Memorize it.
\item Write it down on paper and keep it concealed in a secure place.
\end{enumerate}

\bigskip


\subsection*{Question 65}
What is a password manager?

\begin{enumerate}[label=\Alph*., leftmargin=*]
\item A backup tool for personal identity data.
\item A secure email service.
\item An antivirus software against hacking attacks.
\item A tool that stores, generates, and manages passwords.
\end{enumerate}

\bigskip


\subsection*{Question 66}
What is the main advantage of using a password manager?

\begin{enumerate}[label=\Alph*., leftmargin=*]
\item Remembering only one master password.
\item Eliminating the need for online accounts.
\item Helping to memorize all passwords using mnemonic techniques.
\item No longer creating new passwords.
\end{enumerate}

\bigskip


\subsection*{Question 67}
What are the main features of a password manager?

\begin{enumerate}[label=\Alph*., leftmargin=*]
\item Store, generate, and organize passwords.
\item Backup and organize personal documents.
\item Analyze and remove viruses.
\item Optimize connection speed.
\end{enumerate}

\bigskip


\subsection*{Question 68}
What characteristics define a good password?

\begin{enumerate}[label=\Alph*., leftmargin=*]
\item Short, simple, and easy to remember.
\item Based on personal information.
\item Long, complex, and unique.
\item Identical for all accounts.
\end{enumerate}

\bigskip


\subsection*{Question 69}
Why is it recommended to use randomly generated passwords?

\begin{enumerate}[label=\Alph*., leftmargin=*]
\item To make them easier to remember.
\item To reduce their complexity.
\item To resist brute-force attacks.
\item To resist Sybil attacks.
\end{enumerate}

\bigskip


\subsection*{Question 70}
Why is it essential to secure your email?

\begin{enumerate}[label=\Alph*., leftmargin=*]
\item Because it is used to receive newsletters
\item Because it allows you to reset your passwords
\item Because it contains your banking information
\item Because it contains all of your passwords
\end{enumerate}

\bigskip


\subsection*{Question 71}
What is a fundamental criterion when creating a new email account?

\begin{enumerate}[label=\Alph*., leftmargin=*]
\item Use a unique and strong password.
\item Choose a simple password for easy memorization.
\item Reuse the password from other accounts.
\item Choose a password with the name of your pet.
\end{enumerate}

\bigskip


\subsection*{Question 72}
What advantage does a secure email provider like ProtonMail offer?

\begin{enumerate}[label=\Alph*., leftmargin=*]
\item Integration of an algorithm to automatically generate your emails.
\item Privacy protection and enhanced security.
\item Reduction in features to minimize attack surface.
\item Simpler interface for beginners.
\end{enumerate}

\bigskip


\subsection*{Question 73}
Why is it advised to create two separate email addresses?

\begin{enumerate}[label=\Alph*., leftmargin=*]
\item To receive more spam.
\item To create multiple accounts and increase online visibility.
\item To separate professional and personal communications.
\item To separate general communications from account recovery.
\end{enumerate}

\bigskip


\subsection*{Question 74}
What does the Have I Been Pwned website allow you to do?

\begin{enumerate}[label=\Alph*., leftmargin=*]
\item Create a secure email inbox.
\item Generate new, strong, and random passwords.
\item Check if an email address has been compromised.
\item Synchronize email accounts for easier use.
\end{enumerate}

\bigskip


\newpage
\section*{Answer Key}

\begin{enumerate}
\item \textbf{Question 1:} B
\\ \textit{Private browsing is a feature that prevents internet
browsers from storing browsing history, temporary internet
files, form data, cookies, and user names and passwords.
However, it does not hide browsing from your internet
service provider. On the other hand, TOR network offers
anonymity by masking the user's IP address and allowing
access to the Darknet.}

\item \textbf{Question 2:} B
\\ \textit{While browser extensions can provide useful
functionalities, they can also pose security risks as they
can access your data and browsing activity. Additionally,
having too many extensions can slow down your browser and
affect its performance.}

\item \textbf{Question 3:} A
\\ \textit{Downloading software from the official website of the
publisher ensures that you are getting the legitimate
version of the software, which is free from any malicious
software that could harm your device or steal your data.}

\item \textbf{Question 4:} C
\\ \textit{Open-source software is software whose code is known and
accessible to everyone. This allows for verification, among
other things, that there is no hidden access to steal your
personal data. However, it is not always free, safer, or
more user-friendly than closed-source software.}

\item \textbf{Question 5:} C
\\ \textit{Cookies are files created by websites to store information
on your device. While some sites require these cookies to
function properly, they can also be exploited by
third-party sites, especially for advertising tracking
purposes. Deleting cookies after each visit to a site can
help prevent this.}

\item \textbf{Question 6:} B
\\ \textit{Google Chrome browser is known for its trackers, which can
collect data about your online activity. This can be a
privacy concern for some users.}

\item \textbf{Question 7:} A
\\ \textit{In the context of online browsing, cookies are files
created by websites to store information on your device.
They are used to remember your preferences, login
information, and other data to improve your browsing
experience.}

\item \textbf{Question 8:} A
\\ \textit{Private browsing is a feature that prevents internet
browsers from storing browsing history, temporary internet
files, form data, cookies, and user names and passwords.
However, it does not hide your online activity from your
internet service provider or allow you to browse the
internet anonymously.}

\item \textbf{Question 9:} D
\\ \textit{The TOR (The Onion Router) network is a network that offers
anonymity by masking the user's IP address and allowing
access to the Darknet. It is used by journalists, freedom
activists, and others wishing to escape censorship in
authoritarian countries.}

\item \textbf{Question 10:} A
\\ \textit{In a professional context, VPNs are primarily used to
provide secure remote access to the company's internal
network. This is achieved by encrypting the data exchanged
between the employee and the company's network, making it
difficult for third parties to intercept.}

\item \textbf{Question 11:} D
\\ \textit{Enterprise VPNs are designed for businesses and are
typically more expensive and complex to set up and manage.
On the other hand, consumer VPNs are intended for
individual users and are generally more accessible and
user-friendly.}

\item \textbf{Question 12:} A
\\ \textit{The HTTPS protocol encrypts data on websites,
ensuring that the data exchanged between the user and the
website is secure. This helps protect against third parties
intercepting the data.}

\item \textbf{Question 13:} C
\\ \textit{Verifying the site's public key using a certificate system
is important to confirm the site's authenticity. This
helps protect against malicious individuals impersonating
the site and transferring data in plain text.}

\item \textbf{Question 14:} D
\\ \textit{In the European Union, data protection is regulated by the
General Data Protection Regulation (GDPR). Therefore,
European Wi-Fi access point providers, such as SNCF, do not
resell user connection data, making them safer for personal privacy.}

\item \textbf{Question 15:} D
\\ \textit{Verifying the identity of the site you are browsing is crucial to avoid online scams. You can do this by checking the extension and domain name of the site.}

\item \textbf{Question 16:} B
\\ \textit{When choosing a VPN, it is important to prioritize
reliability and technicality over popularity. This is
because VPN providers that collect the least personal
information are generally the safest, providing better
protection for your online security.}

\item \textbf{Question 17:} B
\\ \textit{Using a VPN at home can enhance the security of the data you exchange online. VPNs encrypt data, making it more difficult for third parties to intercept.}

\item \textbf{Question 18:} B
\\ \textit{Checking the URLs and the small padlock in the address bar
is one way to confirm that you are on the site you intend
to visit. The padlock indicates that the site is using the
HTTPS protocol, which encrypts data for secure exchange.}

\item \textbf{Question 19:} D
\\ \textit{Using an administrator account for daily use on Windows can expose the system to more security risks. Two separate accounts should be created: one for administrative tasks and another for daily use.}

\item \textbf{Question 20:} B
\\ \textit{Regular updates are essential to protect the system from
new threats and vulnerabilities. They are not intended to
make the machine obsolete or slow.}

\item \textbf{Question 21:} C
\\ \textit{Windows Defender is a built-in antivirus program for Windows. It is a safe and effective solution for protecting the system from threats.}

\item \textbf{Question 22:} D
\\ \textit{Cracked software that cannot be updated represents a double potential threat. It can bring a virus during its illegal download from a suspicious website, and be vulnerable to new forms of attack.}

\item \textbf{Question 23:} C
\\ \textit{Inserting unknown or suspicious USB keys can expose the
system to malicious programs that can automatically launch
upon insertion. Checking the USB key will be useless once
it has been inserted.}

\item \textbf{Question 24:} A
\\ \textit{A strong password is crucial for computer security to
prevent unauthorized access to the system.}

\item \textbf{Question 25:} D
\\ \textit{Different operating systems have various solutions for data encryption. BitLocker is for Windows, LUKS is for Linux, and Mac has a built-in option for data encryption.}

\item \textbf{Question 26:} C
\\ \textit{It is recommended that you write down the password for data encryption on paper and keep it in a safe place. This will prevent loss of access to encrypted data.}

\item \textbf{Question 27:} D
\\ \textit{Linux and Mac, thanks to their user rights separation
system, often do not need antivirus.}

\item \textbf{Question 28:} D
\\ \textit{Paying the ransom does not guarantee that you will get your data back. Moreover, it encourages cybercriminals to continue their malicious activities as they see it as a profitable venture.}

\item \textbf{Question 29:} B
\\ \textit{Emails are the first vector of attack for hackers. They use
phishing attempts to extract sensitive information such as
your credentials and passwords.}

\item \textbf{Question 30:} C
\\ \textit{Regularly backing up your data allows you to recover your data in the event of a cyber attack or hardware failure without losing crucial information.}

\item \textbf{Question 31:} B
\\ \textit{The golden rule to protect yourself from fraudulent emails
is to carefully check the full name of the sender as well
as the origin of the email. When in doubt, delete it!}

\item \textbf{Question 32:} A
\\ \textit{If you detect a virus on your computer, disconnect it from the internet to prevent it from spreading or communicating with the hacker. Then, perform a full antivirus scan and delete infected files.}

\item \textbf{Question 33:} A
\\ \textit{Ransomware is a type of malicious software that encrypts user data and demands a ransom to decrypt it. It is becoming increasingly common and can be very troublesome for a company or an individual.}

\item \textbf{Question 34:} A
\\ \textit{While cloud backup services provide a convenient way to
back up your data, they also mean that your data is
potentially on the internet and held by a trusted third
party. This could potentially expose your data to
additional risks.}

\item \textbf{Question 35:} D
\\ \textit{Email attachments and images can contain malware.
Therefore, you should not download or open attachments from
unknown or suspicious senders.}

\item \textbf{Question 36:} B
\\ \textit{Phishing attempts aim to extract sensitive information such as your credentials and passwords. Hackers use various tactics, such as fraudulent emails, to trick you into revealing your information.}

\item \textbf{Question 37:} B
\\ \textit{Self-hosted password managers like Bitwarden or KeePass allow you to control where your passwords are stored, potentially offering more security and control. However, they may require more technical knowledge to use.}

\item \textbf{Question 38:} D
\\ \textit{A password manager generates stronger and longer passwords, offering more security. By gradually replacing your old passwords with these, you can enhance the security of your online accounts.}

\item \textbf{Question 39:} D
\\ \textit{A password manager helps you create and remember strong passwords so you don't have to. It stores all your passwords in a secure way and generates strong passwords for you.}

\item \textbf{Question 40:} A
\\ \textit{A self-hosted password manager requires more technical
knowledge to use, as you need to set up and manage your own
server. However, it can offer more control and security.}

\item \textbf{Question 41:} C
\\ \textit{2FA adds an extra layer of security to your online
accounts. Even if a hacker obtains your password, they will
not be able to access your account without the second
verification factor.}

\item \textbf{Question 42:} B
\\ \textit{The second step in 2FA can be a temporary code sent via SMS, email, a code generated by an application, or a physical security key. However, not all of these methods allow us to keep the same level of security.}

\item \textbf{Question 43:} D
\\ \textit{SMS is not considered the best option for 2FA because it is a 2FA method that is vulnerable to attacks, like the SIM swap scam. Once successfully exploited, these vulnerabilities destroy any protection given by the additional layer of 2FA authentication.}

\item \textbf{Question 44:} D
\\ \textit{One-time passwords (HOTP and TOTP) are safer than SMS
because they require cryptographic calculation and are
stored locally rather than in memory.}

\item \textbf{Question 45:} D
\\ \textit{Cybersecurity is a moving target because threats and
security measures are constantly evolving. It's important
to keep learning and adapting your practices.}

\item \textbf{Question 46:} B
\\ \textit{Hardware tokens offer optimal security because they
generate a unique private key for each site and verify the
URL before allowing the connection. This makes it much
harder for a hacker to gain unauthorized access.}

\item \textbf{Question 47:} B
\\ \textit{Biometric data is less secure because it should remain on
the authentication device and not be disclosed online. If
it is disclosed online, it can be stolen and used for
unauthorized access.}

\item \textbf{Question 48:} A
\\ \textit{For optimal security, it is recommended to use a secure
email address, a secure password manager, and adopt 2FA
using YubiKeys. It is also advisable to purchase two
YubiKeys to anticipate loss or theft.}

\item \textbf{Question 49:} D
\\ \textit{The chapter mentions that a solid education in computer
science, information systems, or a related field is often
the ideal starting point for a career in cybersecurity from
an academic perspective. systems, or a related field, is
often the ideal starting point for a career in
cybersecurity.}

\item \textbf{Question 50:} C
\\ \textit{The chapter emphasizes the importance of practical
experience as it allows you to apply your theoretical
knowledge in real-world situations.}

\item \textbf{Question 51:} D
\\ \textit{The chapter mentions that attending cybersecurity conferences and workshops allows you to learn and helps you build connections with industry professionals.}

\item \textbf{Question 52:} C
\\ \textit{The chapter lists Wireshark, Metasploit, and Nmap as some
of the tools used in cybersecurity.}

\item \textbf{Question 53:} B
\\ \textit{The chapter mentions log analysis and incident response as
some of the skills needed in cybersecurity.}

\item \textbf{Question 54:} A
\\ \textit{The chapter mentions that the constant evolution of threats
in cybersecurity requires regular monitoring.}

\item \textbf{Question 55:} D
\\ \textit{The chapter lists Python, C, and Java as some of the
programming languages used in cybersecurity.}

\item \textbf{Question 56:} B
\\ \textit{The chapter mentions TCP/IP, VPN, and firewall as some aspects of network knowledge needed in cybersecurity.}

\item \textbf{Question 57:} A
\\ \textit{The CISSP certification is one of the most internationally recognized in the field of cybersecurity.}

\item \textbf{Question 58:} A
\\ \textit{Browser extensions that automatically fill in passwords can make the user more vulnerable to phishing attacks. They also tend to increase the attack surface and can slow down browser performance, thus presenting a significant risk.}

\item \textbf{Question 59:} C
\\ \textit{Passkeys offer encrypted random keys combined with a local factor (biometrics or PIN) that a provider hosts. However, the keys remain out of their reach. This approach eliminates the need for passwords, thus providing a high level of security without the constraints associated with traditional passwords.}

\item \textbf{Question 60:} B
\\ \textit{The Single Sign-On (SSO) offered by internet giants poses
problems in terms of availability and risks of censorship.
This is because if the SSO provider experiences an outage
or decides to block a user, the user could lose access to
all services that use that SSO for authentication.}

\item \textbf{Question 61:} C
\\ \textit{Security vulnerabilities are often the result of human
negligence, such as using default passwords or writing down
passwords in insecure places. Sophisticated attacks are not
always necessary to jeopardize computer security.}

\item \textbf{Question 62:} D
\\ \textit{DOH (DNS over HTTPS) and DOT (DNS over TLS) encrypt the DNS connection, providing greater privacy and security. Due to their enhanced security, these protocols are widely used in enterprises and are natively supported by Windows, Android, and iPhone. They improve privacy and security by avoiding man-in-the-middle attacks.}

\item \textbf{Question 63:} C
\\ \textit{Lightning authentication generates a different identifier for each service, without providing an email address or personal information. This can enhance privacy and security by reducing the amount of personal information shared with each service.}

\item \textbf{Question 64:} D
\\ \textit{It is recommended that the master password of a password manager be written down on paper and kept in a concealed and secure location. In fact, if the hard drive is encrypted and the computer is locked, the password will not be accessible, even in the case of burglary.}

\item \textbf{Question 65:} D
\\ \textit{A password manager centralizes the storage and creation of passwords.}

\item \textbf{Question 66:} A
\\ \textit{The user only needs to remember one password to access all the others.}

\item \textbf{Question 67:} A
\\ \textit{The essential functions include storing, generating, and organizing passwords.}

\item \textbf{Question 68:} C
\\ \textit{A good password should be sufficiently long, complex, and should not be reused.}

\item \textbf{Question 69:} C
\\ \textit{Randomness makes passwords harder to guess through trial and error, thereby strengthening their security against brute-force attacks.}

\item \textbf{Question 70:} B
\\ \textit{The mailbox is the central access point for retrieving passwords and compromising multiple accounts in case of a breach.}

\item \textbf{Question 71:} A
\\ \textit{A unique and strong password is essential to prevent hackers from accessing the email account.}

\item \textbf{Question 72:} B
\\ \textit{ProtonMail stands out for its privacy policy and strengthened security measures.}

\item \textbf{Question 73:} D
\\ \textit{Separating email addresses helps limit risks and better protect access to critical accounts.}

\item \textbf{Question 74:} C
\\ \textit{This site helps detect data breaches and verify the security of an email address.}


\end{enumerate}

\end{document}
